% !TEX root = ../main.tex
\section{논의와 한계}
\label{sec:discussion}

\textbf{포스트컨디션을 쓰는 주체.} 사용자가 포스트컨디션을 직접 쓰는 일은 드물다. 애플리케이션의 프런트엔드가 만들고, 지갑이 승인받으려고 보여 준다. ERC-20 승인도 같은 방식으로, 흔히 무제한으로 만들어지므로, 의도의 증거로서 포스트컨디션이 EndorseRank가 쓴 allowance보다 약하지는 않다. 악의적인 프런트엔드가 한도를 느슨하게 잡을 수는 있지만, \PCER는 승인을 실제 유출량으로 평가하므로 느슨한 한도는 아무것도 보태지 않는다. Allow 모드가 흔한 것으로 드러나면(RQ1) 신호의 상당 부분이 포괄 승인이 된다. 포괄 승인을 뺀 민감도 분석이 순위가 그것에 얼마나 기대는지 보여 준다.

\textbf{신원 대신 수수료.} 지갑을 수수료로 가중하면 순위는 활발한 지갑이 쓰는 컨트랙트에 유리해지고, 많이 쓰는 사용자가 가끔 쓰는 사용자보다 크게 셈해진다. 그 대안인 지갑당 한 표가 바로 시빌 지갑이 파고드는 지점이므로(명제~\ref{prop:bipartite}) 이 점을 받아들인다. 약점은 두 가지가 남는다. Stacks 채굴자는 수수료를 받으므로 수수료를 자기 자신에게 낼 수 있다. 상한 $f_{\max}$는 그런 트랜잭션 하나가 살 수 있는 양을 제한하지만 유인을 없애지는 못한다. 또 Stacks의 수수료는 낮아서, 사용자가 적은 범주에서는 공격 비용이 균등 재시작일 때의 몇 배가 되더라도 절대 금액은 작을 수 있다. 다른 가중치인 자본-일은 같은 순간에 여러 지갑에서 재사용할 수 없지만 부유한 지갑에 더 유리하다.

\textbf{가격.} 유출량에 값을 매기려면 가격이 필요하다. 가격은 몫 배분 $s_u(c)$에만 들어가고 가중치 $\omega(u)$에는 들어가지 않는다. 따라서 조작된 가격은 그 토큰을 보낸 지갑들의 가중치를, 그 토큰을 보낸 컨트랙트 쪽으로 옮길 수 있을 뿐이다. 유동성이 없는 토큰은 가격이 없으므로 몫 배분에 영향을 주지 못한다. $T_0$이 미리 정해져 있으므로 $T_0$에 맞춘 가격 펌핑이 가장 뻔한 공격이며, 30일 중앙값은 그 펌핑을 오래 유지하는 데 큰 비용이 들게 한다. ERC-20 연구의 관례인 원래 토큰 단위를 쓴 절제 분석이 가격이 얼마나 중요한지 보여 준다.

\textbf{실행 시점 바인딩.} 실행 시점 trait 바인딩을 빼면 인기 있는 라우터가 증폭기로 쓰이는 것을 막을 수 있다. 그 대신 라우터를 통해서만 거래되는 토큰은 그 토큰을 직접 승인한 사용자와, 코드에서 그 토큰을 이름으로 부르는 풀에게서만 점수를 받는다. 토큰은 가장 많이 검색되는 컨트랙트에 속하므로, 평가에서는 바인딩을 더한 절제 분석을 보고한다.

\textbf{정적 추출.} 파서는 리터럴, 상대, 상수 대상을 푼다. 실행 중에 데이터 변수나 맵에 저장된 principal은 풀지 못하고, ExecutorDAO 패턴 외의 허용 목록은 읽지 않는다. 이런 참조는 빠진 간선이 되어, 그 간선이 가리켰을 컨트랙트의 점수를 낮춘다. 거버넌스 트랜잭션에서 상태 기반 참조를 뽑는 일은 후속 과제이다.

\textbf{클론 묶음.} 정규화는 코드는 같지만 수수료율처럼 중요한 매개변수가 다른 컨트랙트를 하나로 합칠 수 있고, 함수 이름을 바꾸고 형식을 고친 복제본은 찾지 못한다. 가장 많이 쓰인 구성원을 대표로 고르면 인기 있는 복제본이 원저자보다 위에 올 수 있다. 그래서 레지스트리는 묶음마다 배포 이력을 보여 주고 가장 먼저 배포된 구성원을 표시한다.

\textbf{데이터 접근과 재현성.} 호스팅된 Stacks Blockchain API는 요청 빈도를 제한하며 일부 연구 환경에서는 막혀 있다. 그래서 평가는 직접 운영하는 노드와 API를 쓰고, 추출 스크립트와 매니페스트를 공개한다. 트랜잭션 수준 데이터에는 지갑 주소가 들어 있다. 공개된 정보이지만 결과는 집계로만 보고하고, 지갑에는 순위를 매기지도 표시하지도 않으며, 레지스트리는 사용자 목록이 아니라 개수만 보여 준다.

\textbf{새 프로젝트에 대한 불이익.} 사용량에 기반한 순위는 새 컨트랙트를 늦게 알아본다. 레지스트리는 쓰이지 않은 컨트랙트도 목록에 올리고 배포일 순으로 정렬한 보기를 제공하며, 점수가 낮다는 이유로 컨트랙트에 부정적인 표시를 붙이지 않는다.

\textbf{ERC-20 연구와의 관계.} \ref{sec:method}절은 그 연구에서 단순 개수가 좋은 성적을 낸 까닭을 하나 제시한다. H2는 \PCER가 채택과의 일치에서 기준선보다 나쁘지 않기만을 요구한다. \PCER가 노리는 이득은 조작의 값을 높이는 것이며, 이것은 H1이 검정한다.
